\label{ch_FunctionSecretSharing}
\subsubsection{(k,n) Threshold Scheme / Secret Sharing Scheme}
\label{sub_threshold_scheme}
Liu (\cite{liu1968combMath}) describes the general problem as follows: 
\begin{quote}
Eleven scientists are working on a secret project. 
They wish to lock up the documents in a cabinet so that the cabinet can be opened if and only if six or more of the scientists are present.
What is the smallest number of locks needed?
What is the smallest number of keys to the locks each scientist must carry?
\end{quote}
\cite{shamir1979share} shows that the minimal solution uses 462 locks and 252 keys per scientist.
Also, these numbers become exponentially worse when the number of scientists increases.

Secret Sharing is the generalization of this problem.
It takes a Secret $D$ and divides it into $n$ pieces $D_1,...,D_n$ such that:
\begin{itemize}
    \item At least $k\leq n$ pieces are necessarry to make $D$ computable
    \item If less than $k$ pieces are present, $D$ cannot be efficiently computed
\end{itemize}
\cite{shamir1979share} calls this a (k, n) threshold scheme.

\cite{shamir1979share} models a simple scheme based on polynomial interpolation.
$D$ is represented by a random degree $k-1$ polynomial $q(x)=a_0+a_1x+...+a_{k-1}x^{k-1}$ with $a_0=D$.
Each part is then represented by $D_1=q(1),...,D_i=q(i),...,D_n=q(n)$.
Therefore, one can only calculate $D$ if she has knowledge of at least $k-1$ of these values.  

\cite{boyle2015function} describes the same phenomenon as a "secret sharing scheme". 
This research too will use this terminology.

\subsubsection{Additive / Linear Secret Sharing}
\label{sub_additive_secret_sharing}
\cite{boyle2015function} describe additive secret sharing as the simplest type of secret sharing.
The secret is an element of an Abelian Group $G$.
Thus, the secret $s$ is $s\in G$.
$s$ is reconstructed by adding all $p$ shares $s=\sum_{i=0}^p s_i$ (in $G$). 
Every subset of $p-1$ shares reveals nothing about the secret.
This implies every share $s_b$ on its own hides the secret $s$.

\subsubsection{Function Secret Sharing}
\label{sub_function_secret_sharing}